\section{Worked Experiments}
\label{sec:experiments}

To ground the pedagogical claims of Section~\ref{sec:pedagogy} in something a reader can
reproduce rather than take on faith, this section specifies five complete circuits, each
built from the mod's own components and each demonstrating a distinct piece of standard
analog-circuits curriculum: a resistive divider, a first-order RC step response, a
second-order RLC step response, a diode rectifier, and a memristor's characteristic pinched
hysteresis loop. None uses the ideal op-amp, since every one of these five topologies is
standard undergraduate material well before an op-amp is introduced. Every circuit was built
directly into the same persistent world used throughout this paper's own development, laid
out as five side-by-side benches on a flattened platform beginning a few blocks east of the
world's spawn point, so a reader with access to that world can walk up to any of them
immediately rather than reconstruct them from a diagram. Because every predicted value below
follows from closed-form analysis of the exact component values specified (Section~\ref{sec:circuit-solver}
gives the governing equations for each element), a reader who instead prefers to build these
circuits independently, anywhere, need only reproduce the topology and preset values given
here.

\subsection{A methodological note on presets}

None of the mod's block entities persist their right-click-selected preset (resistance,
capacitance, waveform kind, and so on) across a world save and reload -- confirmed directly
against the source, no block entity in the mod overrides Minecraft's block-entity NBT
read/write methods. Consequently every component below, freshly placed, loads at its
hardcoded default preset regardless of which value the experiment actually calls for, and
reaching the intended value is itself the first step of performing the experiment, exactly
as a student wiring up a real bench supply must first turn its dial to the desired voltage
rather than finding it pre-set. Each experiment below states the default and the exact
number of right-clicks needed to reach the intended value, using the fixed cycle order each
preset wraps through (Section~\ref{sec:circuit-solver} and the appendix give the full preset
tables). A power supply or function generator additionally remains electrically
disconnected -- an open circuit, not a 0~V source -- until it receives a redstone signal
(Section~\ref{sec:circuit-solver}); each bench includes a lever placed directly on top of its
source block for exactly this purpose, left off by default so that flipping it marks a clean
$t=0$ for whichever transient the experiment observes.

\subsection{Experiment 1: basic voltage divider}
\label{sec:exp-divider}

The simplest circuit in this set: a DC source and two resistors in series, tapped at their
midpoint. Built as a single row -- power supply, resistor $R_1$, a wire tap, resistor $R_2$
-- closed into a loop by a plain wire return path and a ground reference, ten blocks east of
spawn. Both resistors are left at their default 100~$\Omega$ preset and the supply at its
default 5~V preset, so this experiment requires no right-clicks at all beyond flipping the
lever.

\textbf{Procedure.} Flip the lever on top of the power supply. Pin the oscilloscope probe to
the wire block between the two resistors (the divider's midpoint) to read its absolute
voltage relative to the ground block elsewhere in the same network.

\textbf{Expected result.} By the standard voltage-divider relation,
\begin{equation}
V_{\text{tap}} = V_{\text{supply}} \cdot \frac{R_2}{R_1 + R_2} = 5~\text{V} \cdot \frac{100}{100+100} = 2.5~\text{V},
\end{equation}
a steady 2.5~V once the lever is flipped, with no transient of any kind (every element here
is purely resistive). A useful follow-up, left to the reader: right-click $R_2$ twice
(100~$\Omega \to 1000~\Omega \to 10\,000~\Omega$) and confirm the tap voltage rises to
$5 \cdot 10\,000/10\,100 \approx 4.95$~V, reproducing the divider relation at a different
ratio without changing the topology at all.

\subsection{Experiment 2: simple RC step response}
\label{sec:exp-rc}

A power supply, a resistor, and a capacitor in series, with the capacitor's far lead tied to
ground -- probing the capacitor directly therefore reads its absolute voltage, the classic RC
charging curve. Built 26 blocks east of spawn. Right-click the power supply twice
($5 \to 9 \to 12$~V) and the resistor twice ($100 \to 1000 \to 10\,000\,\Omega$); right-click
the capacitor once ($10 \to 100\,\mu\text{F}$). This gives a time constant
$\tau = RC = 10\,000 \cdot 100\times10^{-6} = 1$~s -- 20 game ticks -- chosen specifically to
be slow enough to watch on the oscilloscope's rolling 200-sample (10~s) history window, yet
fast enough that the whole curve, not just its first fraction, fits comfortably inside that
window.

\textbf{Procedure.} Pin the probe to the capacitor \emph{before} flipping the lever, so the
charging transient is captured from $t=0$. Flip the lever and observe.

\textbf{Expected result.} The standard first-order step response,
\begin{equation}
V_C(t) = V_{\text{supply}}\left(1 - e^{-t/\tau}\right),
\end{equation}
predicts $V_C(\tau) = 12(1 - e^{-1}) \approx 7.58$~V at $t=1$~s, and
$V_C(5\tau) = 12(1 - e^{-5}) \approx 11.92$~V (within 0.7\% of fully charged) at $t=5$~s --
both comfortably inside the 10~s history window, so a single screenshot taken any time after
about 5 seconds should show the full exponential rise settled at its plateau.

\subsection{Experiment 3: simple RLC step response}
\label{sec:exp-rlc}

The same idea as Experiment~\ref{sec:exp-rc}, but with an inductor inserted between the
resistor and capacitor, turning a first-order response into a second-order one capable of
ringing. Built 42 blocks east of spawn. Right-click the power supply twice (to 12~V, as
above). Right-click the resistor \emph{three} times, cycling
$100 \to 1000 \to 10\,000 \to 10~\Omega$ -- the fourth and smallest preset, reached by
wrapping back around past the top of the cycle rather than by cycling downward directly.
Right-click the inductor twice ($0.1 \to 1 \to 5$~H) and the capacitor twice
($10 \to 100 \to 1000\,\mu\text{F}$) -- both to their largest available preset, which gives
the lowest achievable natural frequency and therefore the most visually legible ringing given
the fixed 20~Hz solve rate.

\textbf{Expected result.} With $R=10\,\Omega$, $L=5$~H, $C=1000\,\mu\text{F}$, the standard
series-RLC step response gives an undamped natural frequency, damping ratio, and damped
frequency of
\begin{equation}
\omega_0 = \frac{1}{\sqrt{LC}} \approx 14.14~\text{rad/s}, \qquad
\zeta = \frac{R}{2}\sqrt{\frac{C}{L}} \approx 0.0707, \qquad
\omega_d = \omega_0\sqrt{1-\zeta^2} \approx 14.09~\text{rad/s},
\end{equation}
a lightly damped ring at $f_d = \omega_d/2\pi \approx 2.24$~Hz (a period of about 9 ticks).
The capacitor voltage overshoots its 12~V final value, peaking at
\begin{equation}
V_{\text{peak}} = V_{\text{supply}}\left(1 + e^{-\zeta\pi/\sqrt{1-\zeta^2}}\right) \approx 21.6~\text{V}
\end{equation}
at $t = \pi/\omega_d \approx 0.22$~s (about 4--5 ticks after flipping the lever), before
ringing down and settling to 12~V with a 2\%-settling time of
$4/(\zeta\omega_0) \approx 4$~s -- comfortably spanning close to half of the 10~s history
window, so the ringing itself, not just its first peak, is directly visible without needing
to time the screenshot precisely. This combination -- a real, undamped-looking oscillatory
transient rather than one visibly flattened by numerical damping -- is precisely the
motivation given in Section~\ref{sec:circuit-solver} for using trapezoidal integration over
backward Euler for the reactive-element companion models.

\subsection{Experiment 4: half-wave rectifier}
\label{sec:exp-rectifier}

A function generator driving a diode and a load resistor in series, with the resistor's far
lead grounded so it, too, reads an absolute voltage directly. Built 58 blocks east of spawn.
The diode is placed facing \emph{toward} the generator (its anode is the lead facing the
direction it was placed, per Section~\ref{sec:circuit-solver}), so that current flows readily
generator-to-load on the forward half of each cycle and is blocked on the reverse half. No
right-clicks are required at all: the function generator's own default -- a 5~V, 1~Hz sine,
once its lever is flipped -- and the diode's default silicon preset (approximately 0.7~V
forward drop) and the resistor's default 100~$\Omega$ load are exactly the values used.

\textbf{Procedure.} Flip the lever. For the clearest comparison, pin the function generator
itself (giving the undistorted input, since a component's trace is the voltage across its
own two leads) as one channel of the time-domain probe and the load resistor (giving the
rectified output) as a second channel, so both waveforms are visible on the same
oscilloscope at once.

\textbf{Expected result.} On each positive half-cycle, once the instantaneous input exceeds
the diode's forward voltage, the output approximately follows the input minus that drop,
peaking around $5 - 0.7 = 4.3$~V; on each negative half-cycle the diode is reverse-biased and
the output reads approximately 0~V. At 1~Hz, ten such rectified pulses appear across the 10~s
history window -- a textbook half-wave-rectified sine, clearly distinguishable from the
smooth input sine on the companion channel.

\subsection{Experiment 5: memristor pinched hysteresis loop}
\label{sec:exp-memristor}

The most advanced circuit in this set, and the only one using the X-Y oscilloscope probe
(Section~\ref{sec:architecture}) rather than the time-domain one: a function generator
driving a memristor in series with an ammeter, tracing the memristor's current against its
own voltage. A pinched hysteresis loop through the origin -- present at any drive frequency,
collapsing to a single-valued curve only in the DC limit -- is the experimental signature
that originally motivated the memristor's postulation and its subsequent physical
confirmation (Section~\ref{sec:related-work}). Built 74 blocks east of spawn, with a
frequency module (default 0.5~Hz, requiring no clicks) placed against the function
generator's north face to slow its default 1~Hz down to 0.5~Hz; the generator's default 5~V
amplitude and sine waveform are used unchanged. The memristor is left at its default
$q_{\max} = 10^{-3}$~C preset ($R_{on}=100\,\Omega$, $R_{off}=10\,000\,\Omega$ are fixed,
not presets, per Section~\ref{sec:circuit-solver}).

\textbf{Procedure.} Flip the lever. Pin the memristor \emph{first} with the X-Y probe and
the ammeter \emph{second}: the probe's pinning rule (Section~\ref{sec:architecture}) always
assigns the most recently pinned position to the vertical (Y) axis, so pinning in this order
puts the memristor's voltage on X and the ammeter's current on Y.

\textbf{Expected result.} Unlike the previous four circuits, this one is genuinely nonlinear
-- the memristor's own resistance depends on the full time history of the current through
it, which in turn depends on that same resistance -- so no closed-form prediction is offered
here; the qualitative shape is the point of the experiment. The current-voltage trace should
appear as a pinched loop crossing through the origin on every cycle (the defining signature
of a memristor, as opposed to an ordinary resistor's straight line or an ordinary
reactive element's open ellipse), widest where the driving voltage is largest and
progressively narrowing toward the origin. Because this mod implements the original,
unwindowed HP linear-drift model (Section~\ref{sec:related-work}), a rough charge-balance
estimate for these parameters (amplitude, frequency, and $q_{\max}$) suggests the accumulated
charge over each half-cycle approaches a substantial fraction of $q_{\max}$, so the loop may
visibly flatten as it approaches its boundary within a half-cycle rather than tracing a
smooth figure-eight throughout -- itself a direct, hands-on illustration of the boundary-pinning
limitation that Biolek et al.'s windowed model was later designed to address, discussed
already in Section~\ref{sec:related-work} and Section~\ref{sec:limitations}.
